% Source: https://artofproblemsolving.com/community/c3h2450506_general_formula__product
% Stable identity: aops-topic-2450506-q01
For an integer $a$ write $F(a)=(a-1)(a-2)\cdots(a-n)$, a product of $n$ consecutive integers. Then $P=F(a_1)F(a_2)\cdots F(a_i)$.

\textbf{(a)} If $a\ge n+1$, then $a-n-1\ge0$ and
\[(a-1)!=(a-1)(a-2)\cdots(a-n)\cdot(a-n-1)!,\qquad\text{so}\qquad F(a)=\dfrac{(a-1)!}{(a-n-1)!}.\]
Also $\dbinom{a-1}{n}=\dfrac{(a-1)!}{n!\,(a-1-n)!}$, so $F(a)=n!\dbinom{a-1}{n}$. Multiplying over $j=1,\dots,i$ gives both forms of $P$. $\blacksquare$

\textbf{(b)} It is enough to show that $n!$ divides $F(a)$ for every integer $a$; then $(n!)^i$ divides the product $P$.
\begin{itemize}
\item If $a\ge n+1$, then $F(a)=n!\dbinom{a-1}{n}$ by (a), and a binomial coefficient is an integer.
\item If $1\le a\le n$, one of the factors is $a-a=0$, so $F(a)=0$, which is divisible by $n!$.
\item If $a\le0$, all $n$ factors are negative. Put $c=n+1-a\ge n+1$. Then
\[F(a)=(-1)^n(1-a)(2-a)\cdots(n-a)=(-1)^n(c-n)(c-n+1)\cdots(c-1)=(-1)^nF(c),\]
which is divisible by $n!$ by the first case.
\end{itemize}
So $(n!)^i\mid P$ for all integers $a_1,\dots,a_i$. $\blacksquare$

(For arbitrary real $a_j$ there is no simpler closed form: each factor $F(a_j)$ is a polynomial of degree $n$ in $a_j$. The factorial and binomial forms in (a) are the ``general formula'' for integers $a_j>n$.)
